% SPDX-License-Identifier: GPL-3.0-or-later
\documentclass{article}
\usepackage{amsmath,amssymb,array,diagbox,graphicx,booktabs,hyperref,multirow,adjustbox,tabularx,makecell,pdflscape}
\usepackage[table]{xcolor}
\begin{document}
\listoftables
\listoffigures
\section{Tables and figures}

Click a table cell to edit its text or mathematics. Arrow keys and Tab move
between cells. The Tables ribbon inserts, deletes, duplicates and swaps rows
and columns, and configures horizontal and vertical lines. References, Contents,
inserts lists of tables or figures; Ctrl+click an entry to reach its caption.

\begin{table}[htbp]
\centering
\caption{Values and formulas}
\label{tab:values}
\begin{tabular}{|l|c|r|}
\hline
Quantity & Formula & Value \\
\hline
Area & $\pi r^2$ & $12.57$ \\
Sum & $\sum_{n=1}^3 n$ & $6$ \\
Fraction & $\frac{1}{2}$ & $0.5$ \\
\hline
\end{tabular}
\end{table}

Table~\ref{tab:values} has a caption above its content. Captions are editable
directly in the instant preview; the caption dialog moves or removes them.

\begin{figure}[htbp]
\centering
\includegraphics[width=0.2\linewidth]{../icons/pastex.png}
\caption{The PasTeX icon}
\label{fig:icon}
\end{figure}

Figure~\ref{fig:icon} has its caption below the image. Insert accepts PNG,
JPEG and PDF images. Image settings change the file, width and height.

\begin{tabular}{lp{4cm}}
Text & A paragraph cell can contain a longer sentence that wraps onto
several lines, including \textbf{bold text} and inline mathematics $x^2$. \\
Another row & A tabular can appear directly in body text.
\end{tabular}

\section{Rules and rectangular selections}
\begin{tabular}{lrr}
\toprule
Sample & Before & After \\
\midrule
A & 10 & 12 \\
B & 15 & 18 \\
C & 20 & 24 \\
\bottomrule
\end{tabular}

This grid uses booktabs rules instead of repeated horizontal lines. Select a
rectangle from A's first value to B's last value. The border buttons apply to
the selected cells; repeated border actions preserve a single rule.
Use the lines dialog for a visual overview. Remove a border, then undo.

Duplicate the B row, move it up, then delete it. Repeat with a column. Use arrows
to leave the first or last cell. Copy a rectangular selection and paste it into
another selected range. Caption settings change text, placement and label;
renaming a label updates its references. Try editing the image width as a
number with a text-width unit, and edit its caption independently of the image.

\subsection*{Repeated columns and diagonal headers}
The column specification \verb+*{5}{c|}c+ repeats five centered columns with
right borders, followed by one unruled column. Diagonal labels can contain text
or mathematics; their direction and dimensions are optional settings.

{\renewcommand{\arraystretch}{1.2}\setlength{\extrarowheight}{1pt}
\begin{tabular}{*{5}{c|}c}
\diagbox{Period}{Day} & Mon & Tue & Wed & Thu & Fri \\
\hline
Morning & 1 & 2 & 3 & 4 & 5 \\
Afternoon & 6 & 7 & 8 & 9 & 10
\end{tabular}}

\begin{center}
\diagbox[dir=NE,width=6em,height=3em]{$x^2$}{$y$}
\qquad
\diagbox{Day}{Room}{Time}
\end{center}

\section{Merged cells and column formatting}
Select a rectangular range and choose Table, Merge selected cells. Split cell
keeps the content in the first cell and restores the empty cells around it.
\begin{center}
\begin{adjustbox}{max width=\linewidth}
\begin{tabular}{@{}>{\bfseries}lcr@{}}
\toprule
Group & Quantity & Value \\
\midrule
\multirow{2}{*}{Samples} & A & 12 \\
 & B & 18 \\
\addlinespace[4pt]
\multicolumn{2}{l}{Combined total} & 30 \\
\bottomrule
\end{tabular}
\end{adjustbox}
\end{center}
\section{Flexible columns, multiline cells and colors}
The X columns share the remaining width. Edit the two lines inside each
makecell directly; colors are stored in ordinary xcolor table commands.
\begin{tabularx}{\linewidth}{@{}lXX@{}}
\toprule
\rowcolor{blue!10} Name & Description & Notes \\
\midrule
Flexible & This column wraps a longer sentence to fit the available width.
 & \makecell[l]{First line\\Second line} \\
Colored & \cellcolor{green!15}An emphasized cell & Another flexible column. \\
\bottomrule
\end{tabularx}

\section{Saved boxes}
\newsavebox{\samplebox}
\sbox{\samplebox}{\textsf{Reusable text} with $x^2$}
\typeout{PasTeX saved-box example}
The stored content can appear twice: \usebox{\samplebox} and \usebox{\samplebox}.

\begin{landscape}
\section{Landscape content}
The compiled page is rotated by pdflscape. Continuous editing keeps its
contents upright; use the compiled PDF to check page rotation.
\begin{tabular}{lll}\toprule First&Second&Third\\\midrule A&B&C\\\bottomrule\end{tabular}
\end{landscape}
\end{document}
